\section{Tradition and Traditionalism}

A movement called ``traditionalism'' is apparently the ``very next phase'', so it becomes necessary to periodically engage in the distasteful task of polemics and critique. The problem arises because ``traditionalists'' are making the claim that their ``traditionalism'' is based on the thought of Julius Evola and Rene Guenon.

I realize that Evola is more accessible than Guenon since he writes about trends and politics, but it cannot be emphasized strongly enough that understanding Evola depends on understanding Guenon, despite the sharp differences in certain areas between them. Guenon's writings are more difficult and do not lend themselves readily to short Internet commentaries or radio interviews. Moreover, reading Guenon is not an intellectual exercise, but a spiritual exercise that demands a true metanoia in order to be grasped in their fullness.

Now, one of the foundational works of Guenon, which is mandatory reading for anyone working toward the recovery of Tradition, is \textit{The Reign of Quantity and the Signs of the Times}. Lo and behold, Chapter XXXI is titled ``Tradition and Traditionalism'' and it deals directly with the peculiar manifestations we see happening in our midst, viz., the misuse of the term ``traditionalism''. Since this book is so readily available and well-known, it is difficult to avoid the conclusion that it is deliberate rather than accidental. Guenon writes:

\begin{quotex}
... words are applied to things which they do not fit in any way, and sometimes in a sense directly opposed to their normal meaning. This is one of the most obvious symptoms of the intellectual confusion which reigns everywhere in the present world; but it must not be forgotten that \emph{this very confusion is willed by that which lies hidden behind the whole modern deviation}.
\end{quotex}


Although Guenon concedes the good faith of those involved, he adds:

\begin{quotex}
Nevertheless it must also be recognized that such errors of interpretation and involuntary misconceptions serve the purpose of certain plans so well that is it permissible to wonder whether their growing diffusion may not be due to some of the suggestions that dominate the modern mentality, all of which lead ultimately to nothing less than the destruction of all that is tradition in the true sense of the word.
\end{quotex}


``The modern mentality is no more than the product of a vast collective suggestion'', but there may come a time where the ``resulting state of disorder and disequilibrium becomes so apparent that some people cannot fail to become aware it.'' Unfortunately, Guenon explains:

\begin{quotex}
The very idea of tradition has been destroyed to such an extent that those who aspire to recover it no longer know which way to turn, and are only too ready to accept all the false ideas presented to them in its place and under its name. ... People described as Traditionalists ... only have a sort of tendency or aspiration towards tradition without really knowing anything at all about it; this is the measure of the distance dividing the ``traditionalist'' spirit from the truly traditional spirit, for the latter implies a real knowledge, and indeed in a sense it is the same as that knowledge. In short, the ``traditionalist'' is and can be no more than a mere ``seeker'', and that is why he is always in danger of going astray, not being in possession of the principles which alone could provide him with infallible guidance; and his danger is all the greater because he will find in his path, like so many ambushes, all the false ideas set on foot by the power of illusion which has a capital interest in preventing him from reaching the true goal of his search.
\end{quotex}


In \textit{Crisis of the Modern World}, Guenon mentions the diabolical and satanic spirits that are preventing the cohesion of those dedicated to Tradition. One of the tactics of the adversary is the misdirection mentioned above, that is, the falsification of the true meaning of words. It is interesting that the self-designated leaders of something called ``radical traditionalism'' boast about their intentions to ``seduce'' young men with flashy web sites and multimedia broadcasts in order to build a mass movement. We, on the other hand, are only trying to influence just the precious few who are interested in sound doctrine. Think hard, young men, and be prudent before committing your life.

\begin{center} \textbf{Vincit omnia Veritas} \end{center}



\flrightit{Posted on 2010-04-15 by Cologero}

\begin{center}* * *\end{center}

\begin{footnotesize}
\begin{sffamily}
\texttt{K. L. Anderson on 2010-04-16 at 16:22 said:}

I would suggest being prudent about Traditionalism itself, which we think blocks the true Outward Path to Godhood of biological evolution, although as strictly an Inward Path it can be useful.

\url{http://civilizingthebeast.blogspot.com/}


\hfill

\texttt{Cologero on 2010-04-17 at 17:32 said:}

First all, we are not only prudent about Traditionalism, but you should have noticed that we reject the very concept itself as anything meaningful.

Second, what is possibly meant by the path to Godhood of biological evolution? Evolution is non-teleological, hence there is no path, no goal.

Third, it is genetic. What is the gene for Godhood?

This is Ahrimanic influcence gone amuck ... test the spirits before you adopt such a theory.

\hfill

\texttt{K. L. Anderson on 2010-04-18 at 14:05 said:}

I commented on your blog entry because I felt that if you have the courage to affirm Evola in this age (although it seems to be the thing to do now days) you may have the courage to affirm something of my view of things.

You are so certain that evolution is non-teleological. I disagree, it is teleological, science, tradition and religion missed on this one.

Godhood is epiphenomenal ( not transcendent in this sense) and I assume the ``genes'' for Godhood would be epiphenomenal as well.

Test the Spirits? How do you know I have not? Ahriman may not be the entity you think it is. You need to civilize the beast of materialism, not try to kill it, especially if the material world is the vehicle required in our evolution to Godhood.

I am suggesting a revitalized religion, it does not only look back to Hyperboria, rejecting the modern world, yet it still uses the ancient Inward Path that Evola affirmed. The best of both worlds.

\hfill

\texttt{Cologero on 2010-04-19 at 18:28 said:}

It is not a question of whether I have the courage to affirm something of your view of things, but rather the intelligence.

You do have a certain insight, unfortunately it is partial and incomplete, compounded by your inability to express it coherently.

I am not ``certain'' that evolution is non-teleoligical (did I ever use that word?) ... as you mention, science considers it such, so I was assuming the conventional understanding of evolution. If you want to use words in your own private manner, how am I to anticipate that? Furthermore, you claim that, although science and religion got it wrong, you somehow got it right, without specifying why or how.

To summarize your ``view'': the world-process is immanent in biological evolution and its teleological goal is God.

First of all, for a process to have a goal, there must be something transcendent to the process to define the goal and direct the process towards it.

To wit, your very existence is part of the world process, yet you claim to know the ending well in advance. In other words, you claim to understand the ``idea'' of the world process. That is the nub: the idea is transcendent, it is not biological (your point about epiphenomenal genes is nonsensical.)

The world process, actually the World Soul, is not and can never be God or the Absolute. Yet they are related. I will leave you with a short quote from Solovyov which makes the point as clearly as it can be made:

``In the divine order the all is an absolute organism eternally, while according to the law of natural being, the all gradually becomes an absolute organism in time.''

If that is what you mean, it is a brilliant insight.

\hfill

\texttt{K. L. Anderson on 2010-04-20 at 12:01 said:}

I like your quote by Solovyov. But I don't see the ``all'' exactly that way. Do you really think I have an inability to express these very difficult things coherently? I think I make things very clear. Perhaps today's blog entry below will help, perhaps not.

Transforming The Tiger

The Holy Spirit, which activates the material world in conjunction with evolution, wants the entire world to evolve to Godhood (the Spirit is the Will to Godhood after all ) but this does not always happen at the same time. The force of devolution (the force of Satan ) pulls the material world toward the Great Black Holes of a given universe ( in a multiverse of universes it now seems) and it is this devolutionary force that the Spirit ``combats'' against.

Evola used the term ``riding the tiger'' in suggesting how Traditionalists might deal with the decadent modern world, meaning a way to transform destructive processes into inner liberation. This is the Involutionary Inward Path to the Soul Within, but this is not how I apply the term tiger. I affirm the Twofold Path.

The Spirit affirms the Evolutionary Outward Path of material evolution, which is the vehicle required to reach Godhood. The Spirit does not so much ride the tiger as help transform the tiger from devolution to evolution. The danger in defining liberation as only inner liberation, the goal of the Inward Path, is that it blocks the evolution of the material world to Godhood in what we call the Great Spiritual Blockade.

\url{http://civilizingthebeast.blogspot.com/2010/04/transforming-tiger.html}


\hfill

\texttt{Cologero on 2010-04-26 at 22:00 said:}

There is no innovation in metaphysics. Perhaps it can be said better, adapted to different situations, connections made, but as for novelty, no. So what are you trying to express?

First you say it is all biological, natural and material: a process that is necessarily blind and pointless. Then you introduce the ``Holy Spirit'' which changes everything. Either the spirit dominates matter or the reverse. If the reverse, the lesser cannot ``evolve'' into the higher. Out of what? The higher had to be virtual from the beginning. So the ``goal of the Inward Path'' is to develop the Soul Power, the Will, to actualize the Logos in the material process. Matter cannot create the soul out of itself, not organize itself.

As for the all, there are no two ways to conceptualize it: it is Absolute and Infinite, otherwise it is not ``All''.

\hfill

\texttt{K. L. Anderson on 2010-04-27 at 12:47 said:}

To respond in order to your statements:

Spirit IS matter, but the highest epiphenomenal matter, it ``dominates'' matter in that it can also be defined as the Will To Godhood, so there is no need for a duality here. And, again, evolution is not ``blind and pointless.'' The lesser CAN evolve into the higher, just as the child grows to be a man. The Logos is definition or principle only, which is less than the Object God which the Logos defines. It is time to affirm the real thing, not merely the definition of the real thing. Matter CAN create the Spirit and Soul out of itself if the matter is the highest evolved matter, which is Godhood.

Who says ``there is no innovation in metaphysics?'' Perhaps not in the religions that stem from the Perennial Tradition, which you cling to, but they are all of the Involutionary Inward Path. The ``innovation'' in metaphysics come from the Evolutionary Outward Path, which is really a recognition of the sacred direction of evolution.

\hfill

\texttt{Cologero on 2010-04-27 at 22:46 said:}

I'm afraid the discussion has come to and end. This is what I mean by incoherence: ``Spirit is matter'' is meaningless equation, as spirit is conscious and matter is unconscious. Do you even understand what ``epiphenomenon'' means? You can't simply ``define'' it to mean Will; that is something that needs to be demonstrated. The rest is nonsense, I'm sorry to have to put it that way. I'm sure you have an insight somewhere, but I suggest you pay attention to those who have trod the path before you. I'm trying to save you years of a wasted life.

As for your total dependence on matter ... that is a Satanic notion and you have given in to Christ's first temptation in the desert:

``And the tempter coming said to him: If thou be the Son of God, command that these stones be made bread.''

Your stones (matter) cannot be made into the bread of life. Think carefully about the direction your ideas are taking you.

\hfill

\texttt{F.A. Mercuri on 2010-04-29 at 04:48 said:}

The example of a male child growing into a man does not negate the metaphysical teaching that the ``lesser can not generate the greater''. The case of the child reaching maturity (and old age), is not the lesser producing the greater; rather, demonstrating the completion of a cycle or process that already contains within itself the possibilities of that completion. The child and the adult man are really the same thing all along, in their essence: a human male. That which changes in the human maturation process are the accidents, while the essence remains the same, from cradle to grave, and moreover, within that essence are already contained the gamut of all possible changes that the accidents might undergo--so, this lesser is not generating anything greater or beyond what the essence has already determined as a series of possibilities unfolding within the context of the laws governing the corporeal state (primarily time, space, and motion).

Paul wrote that, ``Now this I say Brothers, that flesh and blood cannot inherit the Kingdom of God; neither does corruption inherit incorruption'' I Corinthians 15:50. Again, the essence of the corporeality does not contain the possibilities of ``evolving'' out of it anything other than that which it is, or contains; however, the Western Tradition depicts man as a tripartite being, the corporeal (soma) representing but one of the modalities of the integral human--consequently, the possibilities of ``evolving'' (beyond the human state) are found within the (regenerated) essence of the psychic, and then noetic ``bodies'' of man--it is the Nous, as a ``ray'' of Logos, which lastly bridges the chasm between the transcendent and imminent, when refined by metanoia, that is, the transformation (regeneration) of the Nous.

Paul continues ``We shall not all sleep, but be changed'' (supra 15:51), at the Resurrection. This corporeality alone is not responsible for the production of the ``spiritual body'', but by the former ``putting on'' incorruptibility which ``changes'' it--implying that this incorruptibility comes from a different source, a superhuman source. Grace working in tandem with human aspiration engender this ``Anothen'', or ``Birth from Above''; the human component beginning with diskarsis, apatheia, and reaching the metanoia--all necessary, due to the fact that natural mans' psychic and noetic faculties are in a degenerated, or ``fallen'' condition.

The Hermetic Tradition also sheds light on this problem. The Salt never singularly has any power of its own to effect a self-transmutation--instead, it is by the action of an initially separated and purified Mercury and Sulfur, then acting upon the Salt, that the Salt then ``evolves''.

\hfill

\texttt{K. L. Anderson on 2010-04-29 at 11:15 said:}

It is perhaps good that you have stopped my attempt to communicate with you. You have called me satanic, unintelligent, full of nonsense, etc... You are not very friendly. I suppose your version of religion does not include being a gentleman to those you define as satanic. That's a sociobiological example of the insider dealing with the outsider, but then, you virtually dismiss biology in your elaborate rationalizations of spirituality versus materiality. So be it. We will see what the future brings.

\hfill

\texttt{Cologero on 2010-04-29 at 17:44 said:}

Mr. Anderson,

I haven't said anything negative or insulting about you personally.

I mentioned that equating spirit and matter is incoherent and the idea that the ``lesser can evolve to the higher'' is nonsense.

I mentioned that world views depending solely on matter is a satanic theory. That is hardly the same as calling you ``Satanic''.

As for you personally, I suggested that you may have an insight, but it needs to be devoloped properly. You need to learn to separate your personality from your theories.

I also mentioned Solovyov who devoted an entire book to the relationship between Spirit and the world process. He has thought the thing through quite deeply and I suggest you take his work into account. I am far from dismissing biology; to the contrary, I value the integral personality who can integrate spirit, soul, and body into a coherent whole. Understanding the proper relationship betwen Spirit and biology (matter) is hardly tantamount to dismissing the latter.

Unfotrunately, there is no possibility of discussion with anyone who denies any difference between spirit and matter. That is the point.

If I am the unconscious puppet of sociobiology, then so are you. That is the dead end where your theory leads.

\end{sffamily}
\end{footnotesize}

